\section*{A retained high-degree reserve in the universal GU row cut}

All logarithms are base two.  For a genuine Boolean function on a finite
uniform cube put
\[
 p(S)=\widehat f(S)^2,\qquad
 W_d=\sum_{|S|=d}p(S),\qquad
 I=\sum_{d\ge1}dW_d,
\]
and retain the source-pinned quantities
\[
 M_1=\sum_i\min\{p(\varnothing),p(\{i\})\},
\]
\[
 D_{\ge2}=\sum_i(B_i-\Gamma_i)
 +\sum_i\sum_{|A|\ge2}K_{i,A},
\]
\[
 V_1=\sum_{i\ne j}(p(\{j\})+p(\{i,j\}))
 \log_2{\pi(\{j\})\over p(\{j\})+p(\{i,j\})}.
\]

\paragraph{Frozen candidate.}
The candidate is
\[
 G\le4I+{3\over4}U.                                      \tag{1}
\]
The preceding source-complete row cut proved
\[
 G\le D_{\ge2}+{1\over2}V_1+{2\over\ln2}(M_1+W_2).       \tag{2}
\]
The earlier assembly replaced $M_1+W_2$ by $I$.  The present step keeps
the exact slack in that replacement.

\paragraph{Reserve identity.}
Define
\[
 Q_{\ge3}:=I-M_1-W_2.                                    \tag{3}
\]
Since $M_1\le\sum_i p(\{i\})=W_1$, one has the exact decomposition
\[
 \boxed{
 Q_{\ge3}=(W_1-M_1)+W_2+\sum_{d\ge3}dW_d
 \ge W_2+3W_{\ge3},}                                    \tag{4}
\]
where $W_{\ge3}=\sum_{d\ge3}W_d$.  In particular,
$Q_{\ge3}>0$ whenever the Fourier degree is at least three.  No lower
bound on the size of the first nonzero high-degree coefficient is assumed.

Substituting (3) into (2) gives the strengthened all-dimensional row cut
\[
 \boxed{
 G\le D_{\ge2}+{1\over2}V_1
 +{2\over\ln2}(I-Q_{\ge3}).}                           \tag{5}
\]
This retains the signed term, every KL row of order at least two, the
complete singleton reuse functional, and all normalization conventions.

\paragraph{Strictly smaller open leaf.}
It is therefore sufficient to prove
\[
 \boxed{
 D_{\ge2}+{1\over2}V_1
 \le\left(4-{2\over\ln2}\right)I+{3\over4}U
       +{2\over\ln2}Q_{\ge3}}                          \tag{6}
\]
only for tensor-prime, positive-residual genuine Boolean functions of
Fourier degree at least three.  Indeed, (5) and (6) cancel the reserve
and give (1).  Degree at most two was already excluded by the independent
bound $G\le2I$, and tensor additivity localizes a positive residual to a
tensor-prime factor.

Using only the weaker explicit part of (4), it is enough to prove
\[
 D_{\ge2}+{1\over2}V_1
 \le\left(4-{2\over\ln2}\right)I+{3\over4}U
 +{2\over\ln2}(W_2+3W_{\ge3}).                         \tag{7}
\]
Thus the first nonzero high-degree mass now supplies a certified payment;
the previous formulation assigned it no credit.

\paragraph{Scope.}
Equations (5)--(7) are strict quantitative reductions, not a universal GU
proof.  The reserve may be arbitrarily small, so it does not by itself
control the signed high-order KL aggregate.  The remaining task is to
show that, under the minimal-child inequalities, the unpaid excess above
the right side of the old cut is at most $(2/\ln2)Q_{\ge3}$, without
discarding $\sum_i(B_i-\Gamma_i)$ or reusing influence.

If (6) is proved, then (1) holds universally.  The audited GU induction
gives $U\le16I$, and the separately verified reduction
$H\le2I+2U$ yields $H\le34I$.
